\section{\texttt{src/baselines/itti\_koch.py}}

Riferimento: Itti, Koch \& Niebur, \emph{A Model of Saliency-Based Visual Attention for Rapid
Scene Analysis}, IEEE TPAMI 1998. Il file più articolato del progetto: una funzione pubblica e
quattro funzioni di supporto.

\subsection{L'idea}

\textbf{Un punto è saliente se spicca rispetto a ciò che lo circonda.} Il modello non riconosce
oggetti: misura contrasto locale su tre canali che il sistema visivo umano elabora in parallelo
nelle prime fasi della visione.

\begin{center}
\begin{tabular}{l p{9.4cm}}
\toprule
\rowcolor{tblHead}\textcolor{white}{\textbf{Canale}} & \textcolor{white}{\textbf{Corrispettivo biologico}} \\
\midrule
Intensità & Contrasto di luminanza, la via più antica ed elementare. \\
\rowcolor{tblAlt} Colore & Opponenza rosso/verde e blu/giallo: è il modo in cui la retina codifica
effettivamente il colore, tramite cellule gangliari antagoniste. \\
Orientazione & Filtri di Gabor a 0°, 45°, 90°, 135°: modello classico delle cellule semplici
della corteccia visiva primaria (V1). \\
\bottomrule
\end{tabular}
\end{center}

\subsection{Le funzioni di supporto}

\subsubsection{\texttt{\_gaussian\_pyramid}}

\begin{idebox}
\begin{pycode}
def _gaussian_pyramid(channel, levels):
    pyramid = [channel.astype(np.float32)]
    for _ in range(levels - 1):
        pyramid.append(cv2.pyrDown(pyramid[-1]).astype(np.float32))
    return pyramid
\end{pycode}
\end{idebox}

Restituisce una \emph{lista} di array: \pyin{pyramid[0]} è l'immagine originale,
\pyin{pyramid[1]} è a metà risoluzione, \pyin{pyramid[2]} a un quarto, e così via.

Ogni livello si costruisce da \pyin{pyramid[-1]}, cioè dall'ultimo aggiunto: la riduzione è
cumulativa.

\begin{nota}
\textbf{A cosa serve la piramide.} Permette di guardare la scena ``a più distanze'': un dettaglio
fine e un oggetto grande producono contrasto a scale diverse, e senza la piramide se ne
coglierebbe solo una.

\medskip
\pyin{cv2.pyrDown} non si limita a scartare pixel: applica prima una sfocatura Gaussiana e poi
sottocampiona. Senza la sfocatura si avrebbe aliasing.
\end{nota}

\subsubsection{\texttt{\_resize\_to}}

\begin{idebox}
\begin{pycode}
def _resize_to(feature_map, shape):
    return cv2.resize(feature_map, (shape[1], shape[0]),
                      interpolation=cv2.INTER_LINEAR)
\end{pycode}
\end{idebox}

Wrapper minimale che esiste per un solo motivo: \textbf{isolare in un punto unico} l'inversione
\pyin{(height, width)} $\to$ \pyin{(width, height)}. Lo scambio \pyin{shape[1], shape[0]} compare
così una volta sola nel modulo, invece che in ogni punto in cui serve un resize.

\subsubsection{\texttt{\_center\_surround}}

\begin{idebox}
\begin{pycode}
def _center_surround(pyramid, center_scales=(2, 3), deltas=(3, 4)):
    maps = []
    for c in center_scales:
        for delta in deltas:
            s = c + delta
            if s >= len(pyramid):
                continue
            center = pyramid[c]
            surround = _resize_to(pyramid[s], center.shape[:2])
            maps.append(np.abs(center - surround))
    return maps
\end{pycode}
\end{idebox}

Doppio ciclo su $c \in \{2,3\}$ (scala \emph{center}, il dettaglio) e $\delta \in \{3,4\}$, con
$s = c + \delta$ (scala \emph{surround}, il contesto). Le quattro combinazioni sono
$(2,5), (2,6), (3,6), (3,7)$.

Il \pyin{continue} salta le combinazioni in cui $s$ eccede i livelli disponibili: con
\pyin{pyramid_levels=7} (indici $0$--$6$) la coppia $(3,7)$ viene scartata, restano tre mappe per
canale.

Prima della sottrazione, la scala grossolana viene riportata alla risoluzione di quella fine:
sono array di dimensioni diverse e non si potrebbero sottrarre direttamente. Il valore assoluto
serve perché conta l'\emph{entità} del contrasto, non il suo segno --- chiaro su scuro e scuro su
chiaro sono entrambi salienti.

È l'analogo computazionale dei campi recettivi center-surround dei neuroni retinici.

\subsubsection{\texttt{\_normalize\_map} --- l'operatore \texorpdfstring{$N(\cdot)$}{N()}}

La funzione concettualmente più densa del progetto.

\begin{idebox}
\begin{pycode}
def _normalize_map(feature_map, M=10.0):
    feature_map = feature_map - feature_map.min()
    if feature_map.max() > 0:
        feature_map = feature_map / feature_map.max() * M

    local_max = cv2.dilate(feature_map, np.ones((7, 7), np.uint8))
    is_local_max = (feature_map == local_max) & (feature_map < M - 1e-3)
    peaks = feature_map[is_local_max]
    mean_peak = peaks.mean() if peaks.size > 0 else 0.0

    return feature_map * (M - mean_peak) ** 2
\end{pycode}
\end{idebox}

\paragraph{Righe 1--3: portare la mappa in $[0, M]$.}
Scala lineare in un intervallo noto, con $M = 10$ (valore del paper). La scelta esatta è
arbitraria: conta il \emph{rapporto} fra $M$ e la media dei picchi.

\paragraph{Riga 5: trovare i massimi locali con una dilatazione.}
\pyin{cv2.dilate} sostituisce ogni pixel con il massimo della sua vicinanza $7\times7$. Di
conseguenza, un pixel che dopo la dilatazione \emph{vale ancora quanto prima} è esattamente un
massimo locale. È un trucco elegante: identifica tutti i massimi locali con una singola
operazione morfologica invece che con una scansione esplicita.

\paragraph{Riga 6: escludere il massimo globale.}
La condizione \pyin{feature_map < M - 1e-3} scarta il picco più alto, che vale esattamente $M$
dopo la normalizzazione. Interessa la media degli \emph{altri} picchi: è quella che dice se la
mappa ha un vincitore netto o tanti competitori.

\paragraph{Righe 7--8: la media dei picchi, con guardia.}
\pyin{peaks.size > 0} protegge dal caso in cui non ci siano massimi locali oltre al globale:
\pyin{.mean()} su un array vuoto restituirebbe NaN con un warning.

\paragraph{Riga 10: il fattore di enfasi.}
\[
N(\text{mappa}) = \text{mappa} \times (M - \bar{m})^2
\]
dove $\bar{m}$ è la media dei massimi locali.

\begin{center}
\begin{tabular}{l l l}
\toprule
\rowcolor{tblHead}\textcolor{white}{\textbf{Situazione}} & \textcolor{white}{$\bar{m}$} & \textcolor{white}{\textbf{Effetto}} \\
\midrule
Un solo picco netto & bassa & $(M-\bar m)^2$ grande $\to$ mappa \textbf{amplificata} \\
\rowcolor{tblAlt} Molti picchi simili & alta & $(M-\bar m)^2$ piccolo $\to$ mappa \textbf{soppressa} \\
\bottomrule
\end{tabular}
\end{center}

\begin{nota}
\textbf{Perché conta.} Due mappe possono avere lo stesso valore massimo ma significato opposto:
una con un picco isolato indica un elemento davvero distintivo, una con venti picchi di altezza
simile indica rumore diffuso.

\medskip
Senza $N(\cdot)$ un canale rumoroso dominerebbe la somma finale per pura forza bruta. È il
meccanismo che rende il modello \emph{selettivo} invece che additivo --- ed è anche la parte del
paper originale che viene più spesso semplificata nelle reimplementazioni, inclusa la nostra.
\end{nota}

\subsection{La funzione principale}

\subsubsection{Estrazione dei canali}

\begin{idebox}
\begin{pycode}
img = image.astype(np.float32)
b, g, r = cv2.split(img)

intensity = (r + g + b) / 3.0

eps = 1e-3
rg = (r - g) / (intensity + eps)
by = (b - np.minimum(r, g)) / (intensity + eps)
\end{pycode}
\end{idebox}

\pyin{cv2.split} separa i tre canali. L'ordine di destrutturazione è \pyin{b, g, r} perché
OpenCV usa BGR.

Le formule di opponenza cromatica riproducono il modo in cui la retina umana codifica il colore:
non come tre canali indipendenti, ma come \emph{differenze} tra coppie antagoniste. Il giallo è
approssimato da \pyin{np.minimum(r, g)}, dato che percettivamente è la componente comune a rosso
e verde.

\begin{nota}
\textbf{Perché dividere per l'intensità.} Normalizza rispetto alla luminosità: un rosso acceso in
ombra e lo stesso rosso in piena luce hanno valori RGB molto diversi ma la stessa \emph{tinta}.
Dividendo per \pyin{intensity} il canale misura la tinta indipendentemente dall'illuminazione.

\medskip
\pyin{eps = 1e-3} evita la divisione per zero nelle zone quasi nere, dove il colore è comunque
privo di significato percettivo.
\end{nota}

\subsubsection{Filtri di Gabor}

\begin{idebox}
\begin{pycode}
orientation_pyramids = []
for theta_deg in (0, 45, 90, 135):
    kernel = cv2.getGaborKernel(
        ksize=(15, 15), sigma=4.0, theta=np.radians(theta_deg),
        lambd=10.0, gamma=0.5, psi=0,
    )
    orientation_pyramids.append(
        [cv2.filter2D(level, cv2.CV_32F, kernel) for level in intensity_pyr]
    )
\end{pycode}
\end{idebox}

Un filtro di Gabor è una sinusoide modulata da una Gaussiana: risponde fortemente ai bordi che
hanno l'orientamento e la frequenza spaziale a cui è sintonizzato.

\begin{center}
\begin{tabular}{>{\ttfamily}l l p{6.6cm}}
\toprule
\rowcolor{tblHead}\textcolor{white}{\textbf{Parametro}} & \textcolor{white}{\textbf{Valore}} & \textcolor{white}{\textbf{Significato}} \\
\midrule
ksize & (15,15) & Dimensione del kernel: deve contenere almeno un periodo della sinusoide. \\
\rowcolor{tblAlt} sigma & 4.0 & Ampiezza dell'inviluppo Gaussiano: quanto è esteso il filtro. \\
theta & 0--135° & Orientamento a cui è sensibile. \pyin{np.radians} converte: OpenCV vuole radianti. \\
\rowcolor{tblAlt} lambd & 10.0 & Lunghezza d'onda della sinusoide: la ``spaziatura'' dei bordi cercati. \\
gamma & 0.5 & Rapporto d'aspetto: allunga il filtro, rendendolo più selettivo sull'orientamento. \\
\rowcolor{tblAlt} psi & 0 & Sfasamento. A 0 il filtro è simmetrico. \\
\bottomrule
\end{tabular}
\end{center}

Il filtro si applica a \emph{tutti} i livelli della piramide di intensità: da qui la list
comprehension interna. Il risultato sono 4 piramidi, una per orientamento.

\pyin{cv2.CV_32F} forza l'output in float: la risposta di Gabor può essere negativa e su
\pyin{uint8} verrebbe troncata.

\subsubsection{Conspicuity map e combinazione finale}

\begin{idebox}
\begin{pycode}
target_shape = intensity_maps[0].shape

def combine(maps):
    resized = [_resize_to(m, target_shape) for m in maps]
    return _normalize_map(np.sum(resized, axis=0))

conspicuity_intensity = combine(intensity_maps)
conspicuity_color = combine(color_maps)
conspicuity_orientation = combine(orientation_maps)

saliency = (conspicuity_intensity
            + conspicuity_color
            + conspicuity_orientation) / 3.0
\end{pycode}
\end{idebox}

Le mappe center-surround provengono da scale diverse e hanno quindi \emph{dimensioni diverse}:
\pyin{combine} le riporta tutte a una risoluzione comune prima di sommarle.
\pyin{np.sum(resized, axis=0)} somma lungo l'asse della lista, cioè elemento per elemento.

Notare la \textbf{doppia normalizzazione}: prima su ogni singola mappa center-surround, poi di
nuovo sulla somma. Serve perché il confronto avviene a due livelli --- prima tra le scale dello
stesso canale, poi tra i tre canali.

La media finale è \emph{non pesata}: i pesi ottimali sarebbero un parametro da apprendere, e il
modello originale è deliberatamente privo di apprendimento.

\subsection{Semplificazioni dichiarate}

Rispetto al paper originale:

\begin{enumerate}
  \item Coppie di scale center-surround fisse ($c \in \{2,3\}$, $\delta \in \{3,4\}$) invece di
  un set più ampio.
  \item Operatore $N(\cdot)$ a singolo passaggio invece che iterativo.
\end{enumerate}

Entrambe riducono il costo computazionale mantenendo il comportamento qualitativo.
\textbf{Vanno dichiarate nel README}: una semplificazione documentata è una scelta di progetto,
una non documentata è un errore.

\section{\texttt{src/annotation/store.py}}

Gestisce la lettura e la scrittura delle annotazioni.

\subsection{Il protocollo}

L'annotatore guarda un'immagine e clicca \textbf{in ordine di importanza decrescente}. Minimo 3
click, massimo 6.

\begin{idebox}
\begin{pycode}
MIN_CLICKS = 3   # sotto, non esiste un ranking sensato
MAX_CLICKS = 6   # sopra, si cercano cose da cliccare invece di notarle
\end{pycode}
\end{idebox}

\begin{nota}
\textbf{Un compito, due ground truth.} L'ordine dei click non è un dettaglio accessorio: rende lo
stesso gesto utile per due scopi.
\begin{itemize}
  \item i click come \emph{punti di fissazione} $\to$ metriche NSS, AUC-Judd;
  \item l'ordine come \emph{ranking} di importanza $\to$ Spearman, Kendall, Top-1.
\end{itemize}
Un secondo giro di annotazione per il ranking raddoppierebbe il lavoro degli annotatori, che è la
risorsa più scarsa.
\end{nota}

\subsection{Il formato su disco}

Un file JSON \textbf{per annotatore}, non un file unico:

\begin{idebox}
\begin{Verbatim}[commandchars=\\\{\},fontsize=\footnotesize]
\textcolor{ideFg}{data/annotations/riccardo.json}

\textcolor{ideFg}{\{}
  \textcolor{ideFunc}{"annotator"}\textcolor{ideFg}{: "riccardo",}
  \textcolor{ideFunc}{"created_at"}\textcolor{ideFg}{: "2026-08-26T18:00:00+00:00",}
  \textcolor{ideFunc}{"annotations"}\textcolor{ideFg}{: \{}
    \textcolor{ideFunc}{"img_001.jpg"}\textcolor{ideFg}{: \{}
      \textcolor{ideFunc}{"image_size"}\textcolor{ideFg}{: [480, 640],}
      \textcolor{ideFunc}{"clicks"}\textcolor{ideFg}{: [\{"x": 412, "y": 230, "order": 1\}, ...]}
    \textcolor{ideFg}{\}}
  \textcolor{ideFg}{\}}
\textcolor{ideFg}{\}}
\end{Verbatim}
\end{idebox}

\begin{nota}
\textbf{Perché un file per annotatore.} Più persone possono annotare in parallelo senza conflitti
di scrittura, ed escludere un annotatore a posteriori diventa banale (si cancella il suo file).
Un file unico richiederebbe un meccanismo di lock, o produrrebbe sovrascritture silenziose.

\medskip
\textbf{Perché salvare \texttt{image\_size}.} Senza le dimensioni, le coordinate sono prive di
significato: se l'immagine venisse ridimensionata o sostituita, i click punterebbero altrove
senza che nulla lo segnali.
\end{nota}

\subsection{\texttt{\_\_init\_\_} --- normalizzazione del nome}

\begin{idebox}
\begin{pycode}
if not annotator or not annotator.strip():
    raise ValueError("Il nome dell'annotatore non puo' essere vuoto.")

self.annotator = annotator.strip().lower().replace(" ", "_")
self.annotations_dir = annotations_dir
self.path = os.path.join(annotations_dir, f"{self.annotator}.json")
self._data = self._load()
\end{pycode}
\end{idebox}

La catena \pyin{strip().lower().replace(" ", "_")} evita che ``Riccardo'', ``riccardo'' e
``Riccardo~'' diventino tre annotatori distinti --- errore facile quando si digita il nome a
inizio sessione.

\subsection{\texttt{\_save} --- scrittura atomica}

\begin{idebox}
\begin{pycode}
def _save(self):
    os.makedirs(self.annotations_dir, exist_ok=True)
    tmp_path = self.path + ".tmp"
    with open(tmp_path, "w", encoding="utf-8") as f:
        json.dump(self._data, f, indent=2, ensure_ascii=False)
    os.replace(tmp_path, self.path)
\end{pycode}
\end{idebox}

Si scrive su un file temporaneo e poi lo si rinomina.

\begin{attenzione}
\textbf{Perché non scrivere direttamente sul file finale.} Se il processo venisse interrotto a
metà scrittura (crash, chiusura del terminale), il file resterebbe troncato --- JSON non valido,
\emph{tutte} le annotazioni precedenti perse.

\medskip
\pyin{os.replace} è \textbf{atomico} sulla maggior parte dei filesystem: o il file nuovo
sostituisce il vecchio per intero, o non succede nulla. Con sessioni da decine di minuti, questa
protezione non è teorica.
\end{attenzione}

\pyin{ensure_ascii=False} preserva gli accenti nei nomi invece di trasformarli in sequenze di
escape.

\subsection{\texttt{save\_annotation}}

\begin{idebox}
\begin{pycode}
if not MIN_CLICKS <= len(clicks) <= MAX_CLICKS:
    raise ValueError(
        f"Servono da {MIN_CLICKS} a {MAX_CLICKS} click, ricevuti {len(clicks)}."
    )
self._data["annotations"][image_name] = {
    "image_size": list(image_size),
    "clicks": [
        {"x": int(x), "y": int(y), "order": i + 1}
        for i, (x, y) in enumerate(clicks)
    ],
    "annotated_at": datetime.now(timezone.utc).isoformat(),
}
self._save()
\end{pycode}
\end{idebox}

L'\pyin{order} si deriva dalla \emph{posizione nella lista}: \pyin{enumerate(clicks)} con
\pyin{i + 1} produce 1, 2, 3\ldots Il chiamante non deve tracciarlo, quindi non può sbagliarlo.

Il salvataggio avviene a ogni chiamata, non a fine sessione: è il salvataggio incrementale.

\subsection{\texttt{load\_all\_annotations}}

\begin{idebox}
\begin{pycode}
for filename in sorted(os.listdir(annotations_dir)):
    if not filename.endswith(".json"):
        continue
    ...
    for image_name, record in data.get("annotations", {}).items():
        if image_name not in aggregated:
            aggregated[image_name] = {
                "image_size": record["image_size"],
                "annotators": {},
            }
        aggregated[image_name]["annotators"][annotator] = record["clicks"]
\end{pycode}
\end{idebox}

Inverte l'organizzazione dei dati: da \emph{per annotatore} a \emph{per immagine}. È la forma
utile per costruire la ground truth, che deve unire le opinioni di tutti su una stessa immagine.

Il filtro \pyin{.endswith(".json")} ignora i file \pyin{.tmp} eventualmente rimasti da una
scrittura interrotta.

\section{\texttt{src/annotation/ground\_truth.py}}

Converte i click in ground truth utilizzabile dalle metriche.

\subsection{Due rappresentazioni, non una}

\begin{center}
\begin{tabular}{l l p{5.8cm}}
\toprule
\rowcolor{tblHead}\textcolor{white}{\textbf{Rappresentazione}} & \textcolor{white}{\textbf{Forma}} & \textcolor{white}{\textbf{Metriche che la usano}} \\
\midrule
Mappa di fissazione & binaria, sparsa & NSS, AUC-Judd \\
\rowcolor{tblAlt} Heatmap continua & Gaussiane sommate & CC, SIM, KL \\
\bottomrule
\end{tabular}
\end{center}

Non è ridondanza: NSS chiede ``quanto vale la saliency \emph{esattamente} sui punti guardati?'' e
sfumare i punti falserebbe la domanda; CC confronta due distribuzioni, e la ground truth deve
essere continua quanto la predizione.

È la stessa doppia rappresentazione dei benchmark di riferimento (MIT/Tübingen).

\subsection{\texttt{clicks\_to\_fixation\_map}}

\begin{idebox}
\begin{pycode}
height, width = image_size
fixation_map = np.zeros((height, width), dtype=np.float32)

for click in clicks:
    x = int(np.clip(click["x"], 0, width - 1))
    y = int(np.clip(click["y"], 0, height - 1))
    fixation_map[y, x] = 1.0

return fixation_map
\end{pycode}
\end{idebox}

Il \pyin{np.clip} protegge da coordinate fuori dai bordi: un click sul bordo estremo
dell'interfaccia può produrre un valore pari alla larghezza, che come indice sarebbe fuori range.
Senza clip si avrebbe un \pyin{IndexError} a metà sessione, con perdita del lavoro in corso.

Notare \pyin{fixation_map[y, x]} --- riga prima, colonna dopo.

\subsection{\texttt{clicks\_to\_heatmap}}

\begin{idebox}
\begin{pycode}
sigma = sigma_frac * width
y_grid, x_grid = np.mgrid[0:height, 0:width].astype(np.float32)

for click in clicks:
    squared_distance = (x_grid - click["x"]) ** 2 + (y_grid - click["y"]) ** 2
    heatmap += np.exp(-squared_distance / (2.0 * sigma ** 2))

max_value = heatmap.max()
if max_value > 0:
    heatmap /= max_value
\end{pycode}
\end{idebox}

\pyin{np.mgrid[0:height, 0:width]} costruisce le griglie di coordinate in una riga, ed è calcolata
\emph{fuori} dal ciclo: viene riusata per ogni click invece di essere ricostruita.

Le Gaussiane si \textbf{sommano}. Dove più click sono vicini i contributi si sovrappongono e il
valore cresce: è così che il consenso tra annotatori emerge automaticamente.

\subsubsection{Il valore di \texorpdfstring{$\sigma$}{sigma}}

\begin{idebox}
\begin{pycode}
DEFAULT_SIGMA_FRAC = 0.05
\end{pycode}
\end{idebox}

\begin{nota}
\textbf{Non è un numero arbitrario.} Il 5\% della larghezza approssima l'ampiezza della fovea
--- circa 1 grado di angolo visivo --- alle distanze di osservazione tipiche di uno schermo.

\medskip
L'idea sottostante: quando si guarda un punto non si vede nitidamente \emph{solo} quel pixel. La
visione ad alta risoluzione copre una piccola regione attorno al punto di fissazione, e la
Gaussiana modella esattamente questo.

\medskip
È espresso come frazione e non in pixel, per lo stesso motivo del center prior: coerenza a
risoluzioni diverse. È una scelta da saper giustificare all'orale.
\end{nota}

\subsection{\texttt{aggregate\_annotators\_heatmap}}

\begin{idebox}
\begin{pycode}
all_clicks = []
for clicks in annotators_clicks.values():
    all_clicks.extend(clicks)
return clicks_to_heatmap(all_clicks, image_size, sigma_frac=sigma_frac)
\end{pycode}
\end{idebox}

Tutti i click di tutti gli annotatori confluiscono in \textbf{un'unica nuvola di punti}, da cui
si genera una sola heatmap.

\begin{nota}
\textbf{L'alternativa scartata.} Costruire una heatmap per annotatore e poi farne la media sembra
equivalente, ma non lo è: la normalizzazione individuale (dividere per il massimo) amplificherebbe
il contributo di chi ha fatto 3 click rispetto a chi ne ha fatti 6. Aggregando prima e
normalizzando dopo, ogni click vale uguale.
\end{nota}
